\documentclass[
  12pt,
  openright,
  a4paper,
  twoside,
  english,
  brazil,
]{article}

\usepackage{mathtools}
\usepackage{amsfonts}
\usepackage{amssymb}
\usepackage{amsmath}
\usepackage[utf8]{inputenc}
\usepackage[portuguese]{babel}
\usepackage[T1]{fontenc}

\usepackage[top=2cm, bottom=2cm, left=2cm, right=2cm]{geometry}
\usepackage{enumerate}

\newcommand{\cqd}{\hfill $\square$}

% -----------------------------------------------------------------------------

\title{Geometria Analítica - Lista 01}
\author{Prof. Helaine - 09/04/2026}
\date{}

\begin{document}

\maketitle

\begin{enumerate}[\bfseries  1.]

  \item Desenho o vetor:
    $$
    \vec{s} = 2\vec{u} + (-3\vec{v}) + \frac{\vec{w}}{2}
    $$

  \item
    $$
    3\vec{w} + 2\vec{v} = \frac{1}{2}\vec{v} + \vec{w}
    $$
    $$
    2\vec{w} = \frac{1}{2}\vec{v} - 2\vec{v}
    $$
    $$
    \vec{w} = \frac{1}{4}\vec{v} - \vec{v} = \frac{1}{4}(-2, 4) - (3, -1)
    $$
    $$
    \vec{w} = \left(-\frac{2}{4} - 3, \frac{4}{4} + 1\right)
    $$
    $$
    \vec{w} = \left(\frac{-14}{4}, 2\right) = \left(-\frac{7}{2}, 2\right)
    $$

  \item
    $$
    (-1, 8) = a_1(1, 2) + a_2(4, -2)\\
    $$

    Sistema:

    $$
    \begin{cases}
      -1 = a_1 + 4a_2 \\
      8 = 2a_1 - 2a_2
    \end{cases}\\
    $$

    $$
    \text{I + 2x II}\implies
    \begin{cases}
      a_1 = 3\\
      a_2 = -1
    \end{cases}
    $$

  \item
    $$
    \vec{CD} = D - C = (x, y) - (-2, 4) = (x + 2, y - 4)\\
    $$
    $$
    (x + 2, y - 4) = \frac{1}{2}(3, -1) - \frac{1}{2}(-1, 2) =
    \frac{1}{2}(4, -3)\\
    $$

    Sistema:
    $$
    \begin{cases} x + 2 = 2 \\ y - 4 = -\frac{3}{2}
    \end{cases} \Rightarrow D = \left(0, \frac{5}{2}\right)
    $$

  \item
    $$
    \vec{CD} = D - C = (x, y) - (-2, 4) = (x + 2, y - 4)\\
    $$
    $$
    (x + 2, y - 4) = \frac{1}{2}(3, -1) - \frac{1}{2}(-1, 2) =
    \frac{1}{2}(4, -3)\\
    $$
    Sistema:
    $$
    \begin{cases} x + 2 = 2 \\ y - 4 = -\frac{3}{2}
    \end{cases} \Rightarrow D = \left(0, \frac{5}{2}\right)
    $$

  \item
    $$
    \frac{m+1}{4} = \frac{3}{2} = \frac{1}{2n-1}
    $$
    $$
    \Rightarrow m + 1 = 3 \cdot 2 \Rightarrow m = 5
    $$
    $$
    2n - 1 = \frac{2}{3} \Rightarrow n = \frac{5}{6}
    $$

  \item
    $$
    M = \frac{1}{2}(x_1 + x_2, y_1 + y_2)
    $$

\end{enumerate}

\end{document}
